\documentclass[11pt,letterpaper,reqno]{amsart}
\usepackage{amssymb}
\usepackage{amsmath}
\usepackage{amsthm}
\usepackage{amsfonts}
\usepackage{enumitem}
\usepackage{booktabs}
\usepackage{graphicx}
\usepackage{xcolor}
\usepackage[T1]{fontenc}
\usepackage{doi}
\usepackage{comment}
\addtolength{\hoffset}{-1.5cm}\addtolength{\textwidth}{3cm}
\addtolength{\voffset}{-1cm}\addtolength{\textheight}{2cm}
\usepackage{hyperref}
\hypersetup{pdfstartview={FitH}}
\usepackage{bookmark}
\usepackage[capitalize,noabbrev]{cleveref}
\usepackage{mathtools}
\usepackage{microtype}

\newtheorem{theorem}{Theorem}[section]
\newtheorem{lemma}[theorem]{Lemma}
\newtheorem{proposition}[theorem]{Proposition}
\newtheorem{corollary}[theorem]{Corollary}
\newtheorem{claim}{Claim}
\newtheorem{question}[theorem]{Question}
\newtheorem{problem}[theorem]{Problem}
\newtheorem{conjecture}[theorem]{Conjecture}
\theoremstyle{definition}
\newtheorem{example}[theorem]{Example}
\newtheorem{remark}[theorem]{Remark}
\newtheorem{definition}[theorem]{Definition}

\newtheorem{theoremletter}{Theorem}
\renewcommand{\thetheoremletter}{\Alph{theoremletter}}
\crefname{theoremletter}{Theorem}{Theorems}
\Crefname{theoremletter}{Theorem}{Theorems}

\newcommand{\R}{\mathbb R}
\newcommand{\tr}{\operatorname{tr}}
\newcommand{\ip}[2]{\left\langle #1,#2\right\rangle}
\newcommand{\norm}[1]{\left\lVert #1\right\rVert}
\newcommand{\ones}{\mathbf 1}

\newcommand{\yinchen}[1]{{
    \color{orange}
    \sf $\star\star\star$ Yinchen: [#1]}}
\newcommand{\quanyu}[1]{{
    \color{purple}
    \sf $\circ\circ\circ$ Quanyu: [#1]}}
\newcommand{\shengtong}[1]{{
    \color{teal}
    \sf $\clubsuit\clubsuit\clubsuit$ Shengtong: [#1]}}

% Bold run-in paragraph heading used in proof outlines
\newcommand{\ppar}[1]{\par\addvspace{\medskipamount}\noindent{\bfseries #1.}\enspace\ignorespaces}

\begin{document}

\title{The Positive and Negative Square-Energy Conjecture}

\author[Y.~Liu]{Yinchen Liu}
\address{Institute for Interdisciplinary Information Sciences, Tsinghua University, Beijing 100084, P. R. China}
\email{liuyinch23@mails.tsinghua.edu.cn}

\author[Q.~Tang]{Quanyu Tang}
\address{School of Mathematics and Statistics, Xi'an Jiaotong University, Xi'an 710049, P. R. China}
\email{tangquanyu827@gmail.com}

\author[S.~Zhang]{Shengtong Zhang}
\address{Stanford University, CA, USA}
\email{stzh1555@stanford.edu}



\subjclass[2020]{Primary 05C50; Secondary 15A18, 15A42}
% 05C50  	Graphs and linear algebra (matrices, eigenvalues, etc.)
% 15A18  	Eigenvalues, singular values, and eigenvectors
% 15A42  	Inequalities involving eigenvalues and eigenvectors


\keywords{positive square energy, negative square energy,
doubly nonnegative matrix, semidefinite programming}



\begin{abstract}
Let $s^+(G)$ and $s^-(G)$ denote the sums of the squares of the positive
and negative adjacency eigenvalues of a graph $G$, respectively. We prove
the conjecture of Elphick, Farber, Goldberg, and Wocjan that every
connected graph $G$ on $n$ vertices satisfies
\[
    \min\{s^+(G),s^-(G)\}\ge n-1.
\]
\end{abstract}





\maketitle



\section{Introduction}

All graphs in this paper are finite, simple, and undirected.  Let
$G=(V,E)$ be a graph with $n=|V|$ vertices, $m=|E|$ edges, and adjacency
matrix $A=A(G)$.  Let $\lambda_1,\ldots,\lambda_n$ be the eigenvalues of
$A$, counted with multiplicity.  The \emph{positive square energy} and
\emph{negative square energy} of $G$ are defined by
\[
    s^+(G):=\sum_{\lambda_i>0}\lambda_i^2,
    \qquad
    s^-(G):=\sum_{\lambda_i<0}\lambda_i^2.
\]
Here and below, $X\succeq0$ means that $X$ is positive semidefinite.
If
\[
    A=A_+-A_-,
    \qquad
    A_+,A_-\succeq0,
    \qquad
    A_+A_-=0
\]
is the decomposition of $A$ into its positive and negative spectral
parts, then
\[
    s^+(G)=\tr(A_+^2),
    \qquad
    s^-(G)=\tr(A_-^2).
\]
Since $\tr(A^2)=2m$, we have
\begin{equation}\label{eq:trace}
    s^+(G)+s^-(G)=2m.
\end{equation}

Elphick, Farber, Goldberg, and Wocjan~\cite[Conjecture~1]{EFGW} proposed the following conjecture.

\begin{conjecture}\label{conj:EFGW}
Every connected graph $G$ on $n$ vertices satisfies
\[
\min\{s^+(G),s^-(G)\}\ge n-1.
\]
\end{conjecture}

Since its formulation, Conjecture~\ref{conj:EFGW} has attracted
considerable attention. The original paper~\cite{EFGW} established the conjecture for several
classes of graphs, including regular and bipartite graphs.  Abiad,
de~Lima, Desai, Guo, Hogben, and Madrid~\cite{AbiadEtAl} proved it for several additional classes.  Further work has
examined structural properties and the asymmetry between the two square
energies~\cite{EL}, as well as a finer conjecture for unicyclic graphs~\cite{NZ}. 


Substantial progress was also made in the general case. Akbari, Kumar,
Mohar, and Pragada~\cite{AKMP} proved the first linear bound
\[
    \min\{s^+(G),s^-(G)\}\ge \frac{3n}{4},
\]
while Zhang~\cite{Zh} simultaneously established the bound
\[
    \min\{s^+(G),s^-(G)\}\ge n-\gamma(G),
\]
where $\gamma(G)$ denotes the domination number of $G$. Recent work has
also obtained refinements concerning the positive square
energy~\cite{AKMPZ}. Related chromatic inequalities and semidefinite or
conic formulations appear in~\cite{AL,CS,CSZ,WE}.

Our main result proves Conjecture~\ref{conj:EFGW}.

\begin{theorem}\label{thm:main}
Every connected graph $G$ on $n$ vertices satisfies
\[
    \min\{s^+(G),s^-(G)\}\ge n-1.
\]
\end{theorem}

The bound is sharp.  Indeed, if $G$ is a tree, then $s^+(G)=s^-(G)=m=n-1$.

Throughout the paper, for a connected graph $G$ we set
\[
    q(G):=2|E(G)|-|V(G)|+1.
\]
By~\eqref{eq:trace}, Theorem~\ref{thm:main} is equivalent to the pair of
upper bounds
\[
      \max\{s^+(G),s^-(G)\}\leq q(G).
\]
We prove these upper bounds.




\subsection{Proof overview}
For real symmetric matrices $X$ and $Y$, write
\[
    \ip XY:=\tr(XY),
    \qquad
    \norm X_F:=\sqrt{\ip XX}.
\]
We use $X\ge0$ for entrywise nonnegativity.  Let $\ones$ be the all-ones
vector, set $J:=\ones\ones^{\top}$, and write $X\circ Y$ for the Hadamard
product of $X$ and $Y$.



Our starting point is the semidefinite method, used throughout
\cite{AL,CS,CSZ,Zh}. They observe that square energies are themselves optimal
values of semidefinite programs.  Indeed, for every $X\succeq0$,
\[
    \ip AX=\ip{A_+}{X}-\ip{A_-}{X}
    \le\ip{A_+}{X}
    \le\norm{A_+}_F\norm X_F
    =\sqrt{s^+(G)}\,\norm X_F,
\]
with equality throughout at $X=A_+$, and therefore
\begin{equation}\label{eq:variational}
    s^+(G)=\max_{0\ne X\succeq0}
    \frac{\max\{\ip AX,\,0\}^2}{\norm X_F^2},
    \qquad
    s^-(G)=\max_{0\ne X\succeq0}
    \frac{\max\{-\ip AX,\,0\}^2}{\norm X_F^2}
\end{equation}
(equivalently, $s^\pm(G)=\min_{M\succeq0}\norm{A\pm M}_F^2$, the form
emphasized in \cite{Zh}).  An upper bound on both square energies is
thus the same thing as a \emph{certificate}: a constant $\beta$ with
\[
    \ip{A}{X}^{2}\le\beta\,\norm X_F^2
    \qquad\text{for every }X\succeq0.
\]

\ppar{Earlier certificates}
Ando and Lin~\cite{AL}, proving a conjecture of Wocjan and
Elphick~\cite{WE}, obtained the certificate
$\beta=2m\bigl(1-\tfrac1{\chi(G)}\bigr)$, and Coutinho and
Spier~\cite{CS} sharpened it to
$\beta=2m\bigl(1-\tfrac1{\chi_v(G)}\bigr)$, where
$\chi_v(G)\le\chi(G)$ is the vector chromatic number. Both arguments use the same mechanism, with a parameter
$t\in\{\chi(G),\chi_v(G)\}$.




Indeed, applying the Cauchy--Schwarz inequality to the $2m$ nonzero
entries of the adjacency matrix gives
\[
    \ip{A}{X}^{2}
    =
    \left(\sum_{u,v}A_{uv}X_{uv}\right)^2
    \le
    2m\sum_{u,v}A_{uv}X_{uv}^{2}
    =
    2m\,\ip{A}{X\circ X}.
\]
The parameter $t$ enters through an inequality on the doubly
nonnegative cone. Namely, setting $Z:=X\circ X$, we have $Z\ge0$
entrywise, and the Schur product theorem gives $Z\succeq0$. Thus $Z$
is doubly nonnegative, and
\begin{equation}\label{eq:dnn}
    \ip{A}{Z}
    \le
    \left(1-\frac{1}{t}\right)\ip{J}{Z}
    \qquad
    \text{for every doubly nonnegative matrix }Z.
\end{equation}
Since $\ip{J}{X\circ X}=\norm{X}_F^2$, these two inequalities, together with~\eqref{eq:variational}, yield \(s^\pm(G) \le 2m\left(1-\frac{1}{t}\right)\).
The conic optimization problem underlying~\eqref{eq:dnn} is studied
systematically by Coutinho, Spier, and Zhang~\cite{CSZ}.



% Uniform Cauchy--Schwarz over the $2m$ nonzero entries of $A$ bounds
% \[
%     \ip{A}{X}^{2}\le 2m\,\ip A{X\circ X},
% \]
% and the parameter $t$ enters through a single positivity claim (whose
% proof is discussed in the next paragraph) about the matrix
% $Z=X\circ X$, which satisfies $Z\ge0$ (entrywise nonnegativity, we use $\geq$ to represent entrywise domination) and, by the Schur
% product theorem, $Z\succeq0$, i.e.\ $Z$ is \emph{doubly nonnegative}:
% \begin{equation}\label{eq:dnn}
%     \ip AZ\le\Bigl(1-\frac1t\Bigr)\ip JZ
%     \qquad\text{for every doubly nonnegative }Z.
% \end{equation}
% Since $\ip J{X\circ X}=\norm X_F^2$, the two bounds combine
% with~\eqref{eq:variational} to give
% $s^\pm(G)\le2m\bigl(1-\tfrac1t\bigr)$.  The conic program behind
% \eqref{eq:dnn} is studied in depth by Coutinho, Spier, and
% Zhang~\cite{CSZ}.

\ppar{The idea behind \eqref{eq:dnn}}
Following the original proofs in~\cite{AL,CS}, the key to~\eqref{eq:dnn}
is the construction of a correlation matrix $C$ (i.e., $C\succeq0$ with
$C_{vv}=1$ for all $v\in V$) such that $C_{uv}\le-\frac1{t-1}$ for
every edge $uv\in E$.  For the parameter $t=\chi(G)$, such a $C$ comes
from a proper coloring made geometric: let
$y_1,y_2,\ldots,y_t\in\R^{t-1}$ be the unit vectors of a regular
simplex centered at the origin, so that
$\ip{y_i}{y_j}=-\frac1{t-1}$ for all $i\ne j$, and set
$c_u:=y_{\text{color of }u}$.  The Gram matrix
$C_{uv}:=\ip{c_u}{c_v}$ is then a correlation matrix with
$C_{uv}=-\frac1{t-1}$ for all adjacent $u$ and $v$, because adjacent
vertices receive distinct colors.  (Relaxing proper colorings to
arbitrary unit vectors with $\ip{c_u}{c_v}\le-\frac1{t-1}$ on edges---a \emph{vector coloring}---is what sharpens the parameter to
$t=\chi_v(G)$~\cite{CS}.)


Note that under this construction $(J-C)\circ A\ge\tfrac t{t-1}A$,
since $A\ge0$ and $1-C_{uv}\ge\tfrac t{t-1}$ on edges; hence
$\ip{(J-C)\circ A}{Z}\ge\tfrac t{t-1}\ip AZ$ for every $Z\ge0$.  The
claim \eqref{eq:dnn} is then an instance of an identity valid for
\emph{every} correlation matrix $C$:
\begin{equation}\label{eq:conic}
    J-(J-C)\circ A
    \;=\;
    \underbrace{\vphantom{(J-C)}C}_{\succeq\,0}
    \;+\;
    \underbrace{(J-C)\circ(J-A)}_{\ge\,0},
\end{equation}
where $J-C\ge0$ because $C\succeq0$ with unit diagonal: every
$2\times2$ principal minor of $C$ gives $|C_{uv}|\le1$.  The
right-hand side is a sum of a positive semidefinite and a nonnegative
matrix, so pairing with any doubly nonnegative $Z$ yields
\[
    \frac t{t-1}\ip AZ\le\ip{(J-C)\circ A}{Z}\le\ip JZ,
\]
the first inequality being the estimate noted above; this is exactly
\eqref{eq:dnn}.

% \ppar{From uniform to weighted certificates}
\ppar{A weighted estimate}
Using the language of correlation-matrix, the sequence of inequality $\ip{A}{X}^{2}\le 2m\ip{A}{X\circ X} \leq 2m \frac{t-1}{t}\ip{(J-C)\circ A}{X\circ X} $ can be recovered by a single Cauchy--Schwarz step, with per-edge weights $1-C_{uv}$:\footnote{In the above proof where $t\in \{\chi(G), \chi_v(G)\}$, $1-C_{uv} = \frac{t}{t-1}$, uniform for any $uv\in E(G)$. }
\[
    \ip{A}{X}^{2}
    =\Bigl(2\sum_{uv\in E}X_{uv}\Bigr)^2
    \le\Phi_G(C)\cdot\ip{(J-C)\circ A}{X\circ X}
    \le\Phi_G(C)\,\norm X_F^2,
\]
where the last inequality uses \eqref{eq:conic}, and
\[
    \Phi_G(C):=\sum_{uv\in E(G)}\frac{2}{1-C_{uv}}.
\]
Now our task is to find a correlation matrix $C$ that minimizes $\Phi_G(C)$ down to $q(G)=2m-n+1$ for every connected graph $G$.

\ppar{Bounding the optimum by duality}
Minimizing $\Phi_G$ over correlation matrices is itself a semidefinite
program: the variable $C$ ranges over the convex set
\[
    \mathrm{Corr}_n:=\{C\in\R^{n\times n}:
    C=C^{\top},\ C\succeq0,\ C_{vv}=1\text{ for all }v\},
\]
cut out by linear equations and one semidefinite constraint. We use the
well-known first-order condition for semidefinite programs (see
\cite[\S3]{VB}) to see what a minimizer $C$ must look like:
complementary slackness produces a matrix
$R\succeq0$ supported off-diagonal on $E(G)$, with
\[
    R_{uv}=\frac1{(1-C_{uv})^2}\quad(uv\in E),
    \qquad
    \sum_{uv\in E}\sqrt{R_{uv}}=\frac{\Phi_G(C)}2,
    \qquad
    \ones^{\top}R\ones=\Phi_G(C).
\]
Now, the most important observation is to use homogeneity and guess the following inequality holds for every connected graph $G$ and every doubly nonnegative matrix $M$ supported on the edges:
\[
    4\Bigl(\sum_{uv\in E}\sqrt{M_{uv}}\Bigr)^2
    \le q(G)\cdot\ones^{\top}M\ones.
\]
Then, taking $M = R$, and the last two identities for $R$ give $\Phi_G(C)\le q(G)$, which is exactly what we want.

Once this observation is established, the proof is surprisingly approachable via induction: an induction on cut vertices---both $q$ and Gram
representations split there---with $2$-connected graphs handled by
averaging the induction hypothesis over all one-vertex deletions.

% \ppar{Dispensing with the program}
\ppar{A direct reduction}
Once the inequality is isolated, the semidefinite program can in fact
be removed from the formal proof altogether: the support hypothesis of the inequality can be dropped by a folding trick, and the unrestricted inequality applied to
$M=A_\pm\circ A_\pm$ then yields Theorem~\ref{thm:main} directly.
This is the route taken in Section~\ref{sec:main}; the
correlation-matrix program remains, in this overview, the account of
how the inequality was found.



\section{A graph-supported doubly nonnegative matrix inequality}\label{sec:sparse}

A real symmetric matrix $X$ is called \emph{doubly nonnegative} if
$X\succeq0$ and $X\ge0$. For a textbook treatment of doubly nonnegative matrices, see \cite{BSM}. Whenever a graph $G$ is fixed and
$u\in V(G)$, let $e_u$ denote the corresponding standard basis vector. Every positive semidefinite matrix $M$ admits a Gram
representation $M_{uv}=\ip{z_u}{z_v}$, where the vectors $z_u$ lie in a Euclidean space and
$\norm{\cdot}$ denotes its Euclidean norm. {\color{red} We use $\sum_{uv \in E(G)}$ to denote summing over each edge $uv$ of $G$ once.}




% For a connected graph $G$ set $q(G):=2|E(G)|-|V(G)|+1$, consistently
% with the introduction.  Edge sums are over unordered edges. The crucial ingredient of our proof is the following strengthening of Theorem~\ref{thm:main} on doubly--nonnegative matrices.


For a connected graph $G$, set
\[
    q(G):=2|E(G)|-|V(G)|+1
\]
as above. The following
inequality is the main ingredient of the proof.

\begin{theorem}\label{thm:sparse}
Let $G$ be a simple connected graph and let $M$ be a doubly nonnegative matrix
indexed by $V(G)$. Then
\[
    4\Bigl(\sum_{uv\in E(G)}\sqrt{M_{uv}}\Bigr)^{\!2}
    \le q(G)\cdot\ones^{\top}M\ones.
\]
\end{theorem}
\begin{proof}
For any pair $(G,M)$ as in the theorem, write
\[
    S(G,M):=\sum_{uv\in E(G)}\sqrt{M_{uv}},
    \qquad
    T(M):=\ones^{\top}M\ones,
\]
so that the claim reads $4S(G,M)^2\le q(G)\,T(M)$.  We induct on
$n:=|V(G)|$.  If $n=1$, the result follows from $\ones^{\top}M\ones \geq 0$. If $n=2$, then $G=K_2$,
$q(K_2)=1$, and $M=\begin{psmallmatrix}a&c\\c&b\end{psmallmatrix}$ with
$c\ge0$; positive semidefiniteness gives
$c\le\sqrt{ab}\le\tfrac{a+b}2$, so
$4S(G,M)^2=4c\le a+b+2c=T(M)$.  Now assume $n\ge3$, and the theorem holds for all graphs with at most $(n - 1)$ vertices.

Next, we will proceed in three steps: first we fold all non-edge entries of $M$ onto the diagonal (hence $M_{uv} = 0$ for all non-edges $uv$ for $u\neq v$); then we treat the case where $G$ has a cut vertex by splitting into two smaller induced subgraphs; finally we handle the case where $G-v$ is connected for every vertex $v$ by averaging the induction hypothesis over all vertex deletions.

\medskip
\noindent\textbf{Folding non-edges.}
Set
\[
    N:=M+\sum_{\substack{u\ne v\\ uv\notin E(G)}}
    M_{uv}\,(e_u-e_v)(e_u-e_v)^{\top},
\]
where the sum runs over all unordered pairs of distinct nonadjacent
vertices.  Each added term is
positive semidefinite, cancels the corresponding non-edge entry of
$M$, and moves its mass to the diagonal; hence $N\succeq0$, $N\ge0$,
$N_{uv}=0$ for non-edges $u\ne v$, and $N_{uv}=M_{uv}$ on edges, so
the edge sums of $M$ and $N$ agree.  Moreover
$\ones^{\top}(e_u-e_v)(e_u-e_v)^{\top}\ones=(1-1)^2=0$, so
$\ones^{\top}N\ones=\ones^{\top}M\ones$. Therefore, without loss of generality, we assume $M_{uv} = 0$ for non-edges $u\ne v$.

\medskip
\noindent\textbf{Cut-vertex case.}
Suppose $G-v$ is disconnected for some $v$.  Group its connected components into
two nonempty unions $U,W$ and set
\[
    G_1:=G[U\cup\{v\}],
    \qquad
    G_2:=G[W\cup\{v\}].
\]
Both are connected (each component of $G-v$ sends an edge to $v$), they
share only $v$, their edge sets partition $E(G)$, and counting vertices
and edges gives
\[
    q(G_1)+q(G_2)=q(G).
\]

Fix a Gram representation $M_{uv}=\ip{z_u}{z_v}$, $u,v\in V(G)$.  No
edge goes between $U$ to $W$, so the support hypothesis makes the spans
$\mathcal U:=\operatorname{span}\{z_u:u\in U\}$ and
$\mathcal W:=\operatorname{span}\{z_w:w\in W\}$ orthogonal.  Consider the orthogonal decomposition of $z_v$
\[
    z_v=p_U+p_W+p_0,
    \qquad
    p_U\in\mathcal U,\quad
    p_W\in\mathcal W,\quad
    p_0\perp\mathcal U+\mathcal W.
\]
Let $M_1$ be the Gram matrix of $(z_u)_{u\in U}$ together with $p_U+p_0$
at $v$, and $M_2$ that of $(z_w)_{w\in W}$ together with $p_W$ at $v$.
Both pairs $(G_1,M_1)$ and $(G_2,M_2)$ satisfy the hypotheses of the
theorem, and every edge entry agrees with the corresponding entry of
$M$: for $u\in U$, $\ip{z_u}{p_U+p_0}=\ip{z_u}{z_v}$, and likewise on
the other side.  Hence the edge sums split,
\[
    S(G,M)=S(G_1,M_1)+S(G_2,M_2).
\]
Moreover $T(M)=\bigl\lVert\sum_{u\in V(G)}z_u\bigr\rVert^2$ splits into
the orthogonal parts $\sum_{u\in U}z_u+p_U$, $\sum_{w\in W}z_w+p_W$,
and $p_0$, so
\[
    T(M)=T(M_1)+T(M_2).
\]
By induction, $2S(G_i,M_i)\le\sqrt{q(G_i)\,T(M_i)}$, so the scalar
Cauchy--Schwarz inequality gives
\[
    2S(G,M)\le\sqrt{q(G_1)T(M_1)}+\sqrt{q(G_2)T(M_2)}
    \le\sqrt{q(G)\,T(M)}.
\]

\medskip
\noindent\textbf{Case with no cut vertex.}
Now $G-v$ is connected for every $v$, so the induction hypothesis can
be averaged over all vertex deletions.  Write $m:=|E(G)|$, $q:=q(G)$,
and split $T(M)$ into diagonal and edge mass,
\[
    T(M)=d_0+2w,
    \qquad
    d_0:=\tr M,
    \quad
    w:=\sum_{uv\in E(G)}M_{uv}.
\]
Deleting a vertex $v$ removes from $S(G,M)$ the terms of the edges at
$v$:
\[
    S(G-v,M-v)=S(G,M)-\sigma_v,
    \qquad
    \sigma_v:=\sum_{u\sim v}\sqrt{M_{uv}},
\]
where $M-v$ denotes the submatrix of $M$ obtained by deleting the row and column corresponding to $v$.  Also, $G-v$ has $n-1$
vertices and $m-d(v)$ edges, $d(v)$ being the degree of $v$, so
$q(G-v)=q+1-2d(v)$.  The induction hypothesis for the pair
$(G-v,M-v)$ therefore reads
\[
    2\bigl(S(G,M)-\sigma_v\bigr)\le\sqrt{q(G-v)\,T(M-v)}.
\]
Sum over all $v$: every edge has two endpoints, so
$\sum_v\sigma_v=2S(G,M)$, and Cauchy--Schwarz gives
\[
    2(n-2)\,S(G,M)
    \le\sum_v\sqrt{q(G-v)\,T(M-v)}
    \le\sqrt{\Bigl(\sum_v q(G-v)\Bigr)\Bigl(\sum_v T(M-v)\Bigr)}.
\]
The two sums are elementary counts.  Using $\sum_v d(v)=2m$ and
$4m=2q+2n-2$, we have
\[
    \sum_v q(G-v)= \sum_v (q + 1 - 2d(v)) = n(q+1)-4m=(n-2)(q-1).
\]
Since each diagonal entry of $M$ survives in $n-1$ of the
submatrices $M-v$ while each edge entry survives in $n-2$ of them, we also have
\[
    \sum_v T(M-v)=(n-1)d_0+2(n-2)w=(n-2)\,T(M)+d_0.
\]
Combining the last three displays and dividing by $(n-2)^2$ yields the
\emph{averaged estimate}
\begin{equation}\label{eq:avg}
    4S(G,M)^2\le(q-1)\,T(M)+\frac{q-1}{n-2}\,d_0.
\end{equation}
Cauchy--Schwarz over the $m$ edges gives a second, \emph{flat},
estimate:
\begin{equation}\label{eq:edgeCS}
    4S(G,M)^2\le4m\sum_{uv\in E(G)}M_{uv}=4mw.
\end{equation}
Now we consider the following two cases, depending on the relative sizes of $w$ and $d_0$.
\ppar{Case 1} If $2(n-1)w\le qd_0$, then $4m=2q+2(n-1)$ turns \eqref{eq:edgeCS} into
\[
    4S(G,M)^2\le4mw=2qw+2(n-1)w\le2qw+qd_0=q\,T(M).
\]
\ppar{Case 2} If instead $2(n-1)w>qd_0$, we use the averaged
estimate \eqref{eq:avg}, for which it suffices that
$\frac{q-1}{n-2}\,d_0\le T(M)$.  Set $\beta:=m-n+1\ge0$.  Since
$q-1=(n-2)+2\beta$ and $T(M)=d_0+2w$, this reduces to
\[
    \beta\,d_0\le(n-2)\,w.
\]
Now simplicity gives $2\beta=2m-2n+2\le(n-1)(n-2)$, so
$q(n-2)-2\beta(n-1)=(n-1)(n-2)-2\beta\ge0$.  Multiplying the case
hypothesis $qd_0<2(n-1)w$ by $n-2$,
\[
    2\beta(n-1)\,d_0\le q(n-2)\,d_0<2(n-1)(n-2)\,w,
\]
and dividing by $2(n-1)$ gives $\beta d_0<(n-2)w$, as needed.
Then \eqref{eq:avg} yields
\[
    4S(G,M)^2\le(q-1)\,T(M)+T(M)=q\,T(M),
\]
completing the induction.
\end{proof}

\section{Proof of the main theorem}\label{sec:main}
In this section, we directly deduce Theorem~\ref{thm:main} from Theorem~\ref{thm:sparse}. 

\begin{proof}[Proof of Theorem~\ref{thm:main}]
The case $n=1$ is trivial, so let $n\ge2$; then $G$ has an edge, and
since $\tr A=0$, the matrix $A$ has eigenvalues of both signs, so
$s^+(G)>0$ and $s^-(G)>0$.

Take $M:=A_+\circ A_+$, which satisfies $M\ge0$ and, by the Schur
product theorem, $M\succeq0$.  Its entries are $M_{uv}=(A_+)_{uv}^2$,
so, using $A_+A_-=0$,
\[
    \sum_{uv\in E(G)}\sqrt{M_{uv}}
    =\sum_{uv\in E(G)}\bigl|(A_+)_{uv}\bigr|
    \ge\sum_{uv\in E(G)}(A_+)_{uv}
    =\tfrac12\ip A{A_+}
    =\tfrac12\tr(A_+^2)
    =\tfrac12\,s^+(G),
\]
while
\[
    \ones^{\top}M\ones=\sum_{u,v}(A_+)_{uv}^2=\norm{A_+}_F^2
    =s^+(G).
\]
Theorem~\ref{thm:sparse} therefore gives
\[
    s^+(G)^2
    \le4\Bigl(\sum_{uv\in E(G)}\sqrt{M_{uv}}\Bigr)^{\!2}
    \le q(G)\,s^+(G),
\]
and dividing by $s^+(G)>0$ yields $s^+(G)\le q(G)$.

For the negative part, take $M:=A_-\circ A_-$. Now
$\ip A{A_-}=-\tr(A_-^2)=-s^-(G)$, so, using $|x|\ge-x$,
\[
    \sum_{uv\in E(G)}\sqrt{M_{uv}}
    =\sum_{uv\in E(G)}\bigl|(A_-)_{uv}\bigr|
    \ge-\sum_{uv\in E(G)}(A_-)_{uv}
    =\tfrac12\,s^-(G),
\]
and $\ones^{\top}M\ones=\norm{A_-}_F^2=s^-(G)$ as before, whence
$s^-(G)\le q(G)$.  By~\eqref{eq:trace},
\[
    s^\pm(G)=2m-s^\mp(G)\ge2m-q(G)=n-1.
    \qedhere
\]
\end{proof}

\begin{remark}
In light of the variational
characterization~\eqref{eq:variational}, this argument says that
Theorem~\ref{thm:sparse} relaxes Theorem~\ref{thm:main} from the
Hadamard squares $\{X\circ X:X\succeq0\}$, which compute the square
energies exactly, to the full doubly nonnegative cone---and the
relaxed inequality still holds with the sharp constant $q(G)$.
Equivalently, taking $M=X\circ X$ in Lemma~\ref{lem:fold} proves that $\ip{A}{X}^{2}\le q(G)\norm X_F^2$ for every $X\succeq0$.
\end{remark}










\section{Consequences of the main theorem}\label{sec:consequences}

We record several consequences of Theorem~\ref{thm:main}.  Throughout this
section, $\kappa(G)$ denotes the number of connected components of $G$, and
\[
    c(G):=|E(G)|-|V(G)|+\kappa(G)
\]
denotes its cyclomatic number.

\subsection{Disconnected graphs}
The following corollary proves the disconnected version of the square-energy
conjecture~\cite[Conjecture~2]{EFGW}.

\begin{corollary}\label{cor:components}
Let $G$ be a graph on $n$ vertices with $\kappa(G)$ connected components.
Then
\[
    \min\{s^+(G),s^-(G)\}\ge n-\kappa(G).
\]
\end{corollary}

\begin{proof}
Let $G_1,\ldots,G_{\kappa(G)}$ be the connected components of $G$, and set
$n_i:=|V(G_i)|$.  Since the adjacency matrix of $G$ is block diagonal,
$s^+$ and $s^-$ are additive over connected components.  Therefore
Theorem~\ref{thm:main} gives, for $\sigma\in\{+,-\}$,
\[
    s^\sigma(G)
    =\sum_{i=1}^{\kappa(G)}s^\sigma(G_i)
    \ge\sum_{i=1}^{\kappa(G)}(n_i-1)
    =n-\kappa(G).
\qedhere\]
\end{proof}

\shengtong{
The following discussion is trivial, let's remove it? 

Combining Corollary~\ref{cor:components} with
$s^+(G)+s^-(G)=2m$, we obtain, for $\sigma\in\{+,-\}$,
\[
    n-\kappa(G)
    \le s^\sigma(G)
    \le 2m-n+\kappa(G).
\]
Equivalently, since $c(G)=m-n+\kappa(G)$, we have
\begin{equation}\label{eq:square-spread-ineq1}
    |s^\sigma(G)-m|\le c(G).
\end{equation}

We also note that $s^+(G)=s^-(G)=n-\kappa(G)$ if and only if $G$ is a forest.  Indeed, if $G$ is a forest, then it is
bipartite and $m=n-\kappa(G)$, so both square energies are equal to
$n-\kappa(G)$.  Conversely, the two equalities imply
$m=n-\kappa(G)$, and hence $G$ is a forest.}



\subsection{Asymmetry of the square energies}

Elphick and Linz~\cite{EL} call $s^+(G)-s^-(G)$ the \emph{squared spread} of
$G$. By \eqref{eq:square-spread-ineq1}, for either sign $\sigma$ we have $|s^\sigma(G)-m|\le c(G)$. It follows that
\begin{equation}\label{eq:square-spread-ineq2}
    |s^+(G)-s^-(G)|\le 2c(G).
\end{equation}
The next result strengthens and makes unconditional the upper bound in
\cite[Proposition~4.1]{EL}.

\begin{corollary}\label{cor:absolute-spread}
Every graph $G$ on $n$ vertices satisfies
\[
    |s^+(G)-s^-(G)|\le (n-1)(n-2).
\]
For $n\ge3$, equality holds if and only if $G\cong K_n$.
\end{corollary}

\begin{proof}
Let $\kappa:=\kappa(G)$.  If the connected components of $G$ have orders
$n_1,\ldots,n_\kappa$, then
\[
    |E(G)|
    \le \sum_{i=1}^{\kappa}\binom{n_i}{2}
    \le \binom{n-\kappa+1}{2}.
\]
The second inequality is attained only when one component has
$n-\kappa+1$ vertices and all the remaining components are isolated
vertices.  It follows that
\[
    c(G)
    =|E(G)|-n+\kappa
    \le \binom{n-\kappa}{2}
    \le \binom{n-1}{2}.
\]
By \eqref{eq:square-spread-ineq2}, we obtain
\[
    |s^+(G)-s^-(G)|
    \le 2c(G)\le 2\binom{n-1}{2}
    =(n-1)(n-2).
\]

Suppose that $n\ge3$ and equality holds.  Then equality must hold in
\[
    c(G)\le\binom{n-\kappa}{2}\le\binom{n-1}{2},
\]
which forces $\kappa=1$ and $|E(G)|=\binom n2$.  Hence $G\cong K_n$.
Conversely, the adjacency eigenvalues of $K_n$ are $n-1$ and $-1$ with
multiplicity $n-1$, so
\[
    s^+(K_n)-s^-(K_n)
    =(n-1)^2-(n-1)
    =(n-1)(n-2).
\qedhere\]
\end{proof}





\subsection{Adjacency inertia}

For a graph $G$, let $n^+(G)$, $n^0(G)$, and $n^-(G)$ denote,
respectively, the numbers of positive, zero, and negative eigenvalues of its
adjacency matrix, counted with multiplicity.  Let $\iota(G)$ denote the
number of isolated vertices.

Zhang~\cite[Corollary~1.5]{Zh} proved that every graph without isolated
vertices satisfies
\[
    \min\{s^+(G),s^-(G)\}
    \ge
    \max\{n^+(G),n^0(G),n^-(G)\}.
\]
The following result proves the stronger conjecture of Elphick recorded
in~\cite[Conjecture~7.4]{Zh} and extends it to
graphs with isolated vertices.

\begin{corollary}\label{cor:inertia}
Every graph $G$ satisfies
\[
    \min\{s^+(G),s^-(G)\}
    \ge
    n^0(G)-\iota(G)+\max\{n^+(G),n^-(G)\}.
\]
\end{corollary}

\begin{proof}
Let $\eta:=\kappa(G)-\iota(G)$ denote the number of connected components of $G$ with at least two vertices.  The adjacency matrix
of any such component has at least one positive and at least one negative eigenvalue.
Hence
\[
    \min\{n^+(G),n^-(G)\}\ge \eta.
\]
Since $n=n^+(G)+n^0(G)+n^-(G)$, we obtain
\[
\begin{aligned}
    n^0(G)-\iota(G)+\max\{n^+(G),n^-(G)\}
    &=n-\iota(G)-\min\{n^+(G),n^-(G)\}\\
    &\le n-\iota(G)-\eta\\
    &=n-\kappa(G).
\end{aligned}
\]
The result follows from Corollary~\ref{cor:components}.
\end{proof}



\subsection{Positive and negative \texorpdfstring{$p$}{p}-energies}
For $p>0$, the corresponding positive and negative $p$-energies are
defined by
\[
    \mathcal E_p^+(G)
    :=\sum_{\lambda_i>0}|\lambda_i|^p,
    \qquad
    \mathcal E_p^-(G)
    :=\sum_{\lambda_i<0}|\lambda_i|^p.
\]
These quantities have also been studied in several recent
papers~\cite{AKMPp,CWZ,ETZ,LTPath,TLW}. Our main theorem resolves a number of conjectures surrounding these energies.

For connected graphs, the next result proves
\cite[Conjecture~5.4]{ETZ}.  Its negative-sign case for $p\ge2$ also proves
\cite[Conjecture~2]{AKMPp}.

\begin{corollary}\label{cor:p-energy}
Let $G$ be a graph on $n$ vertices. For every $p>0$ and $\sigma\in\{+,-\}$,
\[
    \mathcal E_p^\sigma(G)\ge
    \begin{cases}
        \left(n-\kappa(G)\right)^{p/2},&0<p\le2,\\[2mm]
        n-\kappa(G),&p\ge2.
    \end{cases}
\]
\end{corollary}

\begin{proof}
Set $N:=n-\kappa(G)$. The case $N=0$ is immediate.  Suppose that $N>0$, fix
$\sigma\in\{+,-\}$, and let $x_1,\ldots,x_r$ be the absolute values of the
adjacency eigenvalues of sign $\sigma$.  By
Corollary~\ref{cor:components},
\[
    \sum_{i=1}^r x_i^2=s^\sigma(G)\ge N.
\]
If $0<p\le2$, then
\[
    \mathcal E_p^\sigma(G)
    =\sum_{i=1}^r x_i^p
    \ge
    \left(\sum_{i=1}^r x_i^2\right)^{p/2}
    \ge N^{p/2}.
\]
Now let $p\ge2$.  By H\"older's inequality,
\[
    \mathcal E_p^\sigma(G)
    \ge
    r^{1-p/2}
    \left(\sum_{i=1}^r x_i^2\right)^{p/2}.
\]
Every nontrivial connected component of $G$ contributes an adjacency
eigenvalue of sign opposite to $\sigma$, while every isolated vertex
contributes a zero eigenvalue.  Thus at least $\kappa(G)$ eigenvalues are
not among $x_1,\ldots,x_r$, and hence $r\le n-\kappa(G)=N$. Since $1-p/2\le0$, it follows that
\[
    \mathcal E_p^\sigma(G)
    \ge
    N^{1-p/2}N^{p/2}
    =N.
\qedhere\]
\end{proof}

For a connected graph, Corollary~\ref{cor:p-energy} gives
\[
    \min\{\mathcal E_p^+(G),\mathcal E_p^-(G)\}
    \ge(n-1)^{p/2}
    \qquad (0<p\le2).
\]
This proves~\cite[Conjecture~5.4]{ETZ}; the range $0<p<1$ was previously
proved in~\cite[Theorem~5.7]{ETZ}. With the convention $\mathcal E_0^\pm(G):=n^\pm(G)$, the corresponding assertion at $p=0$ is immediate.

If $G$ is connected and $p\ge2$, then the negative-sign bound gives
\[
    \mathcal E_p^-(G)\ge n-1=\mathcal E_p^-(K_n).
\]
This proves~\cite[Conjecture~2]{AKMPp}.  The conjecture was previously known
for all real $p\ge4$ by~\cite{AKMPp} and for every integer $p\ge3$
by~\cite{CWZ}.

The bounds in Corollary~\ref{cor:p-energy} are sharp.  For $0<p<2$,
equality is attained by the disjoint union of a star and isolated vertices.
For $p=2$, equality holds for every forest.  For $p>2$, equality is attained
by a disjoint union of edges and isolated vertices.  In the connected
negative-sign case, equality is attained by $K_n$.



For connected graphs and $p>2$, Corollary~\ref{cor:p-energy} gives
\[
    \mathcal E_p^+(G)\ge n-1.
\]
The stronger conjecture
\[
    \mathcal E_p^+(G)\ge\mathcal E_p^+(P_n),
\]
where $P_n$ is the path on $n$ vertices, remains open in general.  It is
known for connected bipartite graphs for every real $p\ge2$, for all
connected graphs when $p\ge3$ is an odd integer, and for all connected
graphs when $p=4$; see~\cite{LTPath}.


\subsection{Graph complements}

Let $\overline G$ denote the complement of $G$.  Elphick and
Aouchiche~\cite[Conjecture~6]{EA} conjectured that
\[
    s^+(G)+s^+(\overline G)\le(n-1)^2
\]
and observed that this inequality would follow from
Conjecture~\ref{conj:EFGW}.  The following corollary proves their conjecture
and its negative-sign analogue.

\begin{corollary}\label{cor:nordhaus-gaddum}
Let $G$ be a graph on $n\ge2$ vertices.  For every
$\sigma\in\{+,-\}$,
\[
    s^\sigma(G)+s^\sigma(\overline G)\le(n-1)^2.
\]
For $\sigma=+$, equality holds if and only if $G$ is complete or edgeless.
For $\sigma=-$, equality holds if and only if $n=2$.
\end{corollary}

\begin{proof}
Fix $\sigma\in\{+,-\}$, and let $\tau$ denote the opposite sign.  Write $\kappa:=\kappa(G)$ and $\overline\kappa:=\kappa(\overline G)$. By Corollary~\ref{cor:components},
\[
    s^\tau(G)+s^\tau(\overline G)
    \ge 2n-\kappa-\overline\kappa.
\]
Since
\[
    s^+(G)+s^-(G)+s^+(\overline G)+s^-(\overline G)
    =n(n-1),
\]
we obtain
\[
    s^\sigma(G)+s^\sigma(\overline G)
    \le n(n-1)-2n+\kappa+\overline\kappa.
\]
At least one of $G$ and $\overline G$ is connected, and hence $\kappa+\overline\kappa\le n+1$. Therefore
\[
    s^\sigma(G)+s^\sigma(\overline G)\le(n-1)^2.
\]

Suppose that equality holds. Then $\kappa+\overline\kappa=n+1$. Since at least one of $G$ and $\overline G$ is connected, the other has
$n$ connected components.  Thus one of the two graphs is complete and the
other is edgeless.  Now
\[
    s^+(K_n)=(n-1)^2,
    \qquad
    s^-(K_n)=n-1,
\]
whereas both square energies of an edgeless graph vanish.  The equality
statements follow.
\end{proof}






\section*{Acknowledgments and AI disclosure}
Q.~Tang thanks Minghua Lin for helpful discussions.


During the development of this work, all core mathematical ideas and claims were generated independently by the authors. The authors used AI (ChatGPT 5.6 Sol Pro) as a preliminary programming tool to evaluate the plausibility of conjectured claims and suggest initial proof strategies. The tool was not treated as a mathematical authority. While AI provided exploratory pathways, the authors substantially modified and simplified these suggested strategies to construct the final arguments. No mathematical claim, computation, or proof was included without rigorous independent verification. The authors finalized all theorem statements, proofs, and references, and take full responsibility for the correctness, originality, and integrity of the paper.



\begin{thebibliography}{99}



\bibitem{AbiadEtAl}
A.~Abiad, L.~de~Lima, D.~N.~Desai, K.~Guo, L.~Hogben, and J.~Madrid.
\newblock Positive and negative square energies of graphs.
\newblock \emph{Electron. J. Linear Algebra}, 39:307--326, 2023.

\bibitem{AKMP}
S.~Akbari, H.~Kumar, B.~Mohar, and S.~Pragada.
\newblock A linear lower bound for the square energy of graphs.
\newblock \emph{Electron. J. Combin.},
32(3):Paper No.~P3.53, 2025.

\bibitem{AKMPp}
S.~Akbari, H.~Kumar, B.~Mohar, and S.~Pragada.
\newblock Vertex partitioning and $p$-energy of graphs.
\newblock \emph{Linear Algebra Appl.}, 724:96--107, 2025.

\bibitem{AKMPZ}
S.~Akbari, H.~Kumar, B.~Mohar, S.~Pragada, and S.~Zhang.
\newblock Refinement of a conjecture on positive square energy of graphs.
\newblock arXiv:2506.07264, 2025.

\bibitem{AL}
T.~Ando and M.~Lin.
\newblock Proof of a conjectured lower bound on the chromatic number of a
graph.
\newblock \emph{Linear Algebra Appl.}, 485:480--484, 2015.

\bibitem{BSM}
A. Berman and N. Shaked-Monderer.
\newblock Completely positive matrices,
\newblock World Sci. Publ., River Edge, NJ, 2003.

\bibitem{CWZ}
Z.~Chen, Z.~Wang, and X.-D.~Zhang.
\newblock Positive and negative $3$-energies of graphs.
\newblock arXiv:2604.15656, 2026.

\bibitem{CS}
G.~Coutinho and T.~J.~Spier.
\newblock Sums of squares of eigenvalues and the vector chromatic number.
\newblock arXiv:2308.04475, 2023.

\bibitem{CSZ}
G.~Coutinho, T.~J.~Spier, and S.~Zhang.
\newblock Conic programming to understand sums of squares of eigenvalues of
graphs.
\newblock arXiv:2411.08184, 2024.

\bibitem{EA}
C.~Elphick and M.~Aouchiche.
\newblock Nordhaus--Gaddum and other bounds for the sum of squares of the
positive eigenvalues of a graph.
\newblock \emph{Linear Algebra Appl.}, 530:150--159, 2017.

\bibitem{EFGW}
C.~Elphick, M.~Farber, F.~Goldberg, and P.~Wocjan.
\newblock Conjectured bounds for the sum of squares of positive eigenvalues
of a graph.
\newblock \emph{Discrete Math.}, 339(9):2215--2223, 2016.

\bibitem{EL}
C.~Elphick and W.~Linz.
\newblock Symmetry and asymmetry between positive and negative square
energies of graphs.
\newblock \emph{Electron. J. Linear Algebra}, 40:418--432, 2024.

\bibitem{ETZ}
C.~Elphick, Q.~Tang, and S.~Zhang.
\newblock A spectral lower bound on chromatic numbers using $p$-energy.
\newblock \emph{European J. Combin.}, 132:104252, 2026.

\bibitem{LTPath}
Y.~Liu and Q.~Tang.
\newblock Path-minimality for positive $p$-energies, Laplacian-type
spectra, and line graphs.
\newblock arXiv:2606.30996v1, 2026.

\bibitem{NZ}
B.~Ning and J.~Zeng.
\newblock A proof of a conjecture on positive and negative square energies
of unicyclic graphs.
\newblock arXiv:2605.24668, 2026.

\bibitem{TLW}
Q.~Tang, Y.~Liu, and W.~Wang.
\newblock On the positive and negative $p$-energies of graphs under edge
addition.
\newblock \emph{Discrete Appl. Math.}, 388:25--33, 2026.

\bibitem{VB}
L.~Vandenberghe and S.~Boyd.
\newblock Semidefinite programming.
\newblock \emph{SIAM Rev.}, 38(1):49--95, 1996.

\bibitem{WE}
P.~Wocjan and C.~Elphick.
\newblock New spectral bounds on the chromatic number encompassing all
eigenvalues of the adjacency matrix.
\newblock \emph{Electron. J. Combin.},
20(3):Paper No.~P39, 2013.

\bibitem{Zh}
S.~Zhang.
\newblock Extremal values for the square energies of graphs.
\newblock arXiv:2409.15504v2, 2024.

\end{thebibliography}

\end{document}
